\documentclass[11pt]{amsart}
\usepackage[T1]{fontenc}
\usepackage[margin=1in]{geometry}
\usepackage{amsmath,amssymb,amsthm}
\usepackage{lmodern,microtype}
\usepackage[hidelinks]{hyperref}
\hypersetup{pdftitle={Pruned combs and the associative spectrum of SC258},
  pdfauthor={Xiaojian Zeng},pdfsubject={Associative spectrum of a three-element groupoid},
  pdfkeywords={associative spectrum, bracketing, labeled comb, normal form}}
\newtheorem{theorem}{Theorem}
\newtheorem{lemma}[theorem]{Lemma}
\newtheorem{proposition}[theorem]{Proposition}
\newcommand{\B}{\mathcal{B}}
\newcommand{\Sset}{\mathcal{S}}
\newcommand{\Cset}{\mathcal{C}}
\newcommand{\SC}{\mathsf{SC}258}
\title[The associative spectrum of SC258]{Pruned combs and the associative spectrum of SC258}
\author{Xiaojian Zeng}
\address{PBC School of Finance, Tsinghua University, 100190}
\email{xiaojianzeng@foxmail.com}
\date{}
\thanks{2020 Mathematics Subject Classification. 05A15, 20N02, 08B05.}
\keywords{Associative spectrum, binary operation, bracketing, labeled tree, normal form}

\begin{document}
\begin{abstract}
We determine the associative spectrum of the three-element groupoid
\(\SC\). If \(a_n\) counts the distinct functions induced by the full
binary bracketings of \(x_1\cdots x_n\) in this fixed order, then
\(a_1=a_2=1\), \(a_3=2\), and
\(a_n=3a_{n-1}-2a_{n-2}+a_{n-3}\) for every \(n\geq4\).
The proof gives a complete equality criterion for the term functions:
recursively discard the right subtree whenever both children are composite,
retaining all surviving leaf labels. Two bracketings induce the same function
exactly when these pruned trees coincide. We reconstruct each pruned comb
from its essential variables and coordinate restrictions, then enumerate the
attainable forms by their largest retained label. This establishes the
all-arity recurrence suggested by the finite data of Huang and Lehtonen.
\end{abstract}
\maketitle

\section{The spectrum and the main result}

The associative spectrum of a binary operation counts the different term
functions obtained by changing parentheses while keeping the variables in
their original order; this notion was introduced by Cs\'ak\'any and
Waldhauser~\cite{CsakanyWaldhauser2000}. We consider the groupoid
\(A=\{0,1,2\}\) with the following operation, called \(\SC\) in
Huang and Lehtonen~\cite[Section~7]{HuangLehtonen2024}:
\begin{equation}\label{eq:table}
\begin{array}{c|ccc}
 *&0&1&2\\\hline
 0&0&0&0\\
 1&1&1&0\\
 2&1&0&1
\end{array}.
\end{equation}
Rows index left operands and columns index right operands.
Let \(\B_n\) be the set of full binary bracketings whose leaves, read
from left to right, are exactly \(x_1,\ldots,x_n\), each used once.
For \(t\in\B_n\), write \(t^A\colon A^n\to A\) for its induced
function, and put
\[
 a_n=\bigl|\{t^A:t\in\B_n\}\bigr|\qquad(n\geq1).
\]
Equality is tested on every input in \(A^n\); the positions of ignored
variables are retained. Allowing permutations of the variables instead
gives the associative-commutative spectrum of Huang and
Lehtonen~\cite{HuangLehtonen2023}, which is outside the present result.

Huang and Lehtonen~\cite[Section~7]{HuangLehtonen2024} reported
\[
 1,1,2,5,12,28,65,151,351
\]
as the first nine values for \(\SC\), observed agreement with a trisection
of the Padovan sequence, and asked for a general determination. The following
theorem settles that fixed-order counting problem.

\newpage
\begin{theorem}\label{thm:main}
For the operation~\eqref{eq:table},
\begin{equation}\label{eq:recurrence}
 a_1=a_2=1,\qquad a_3=2,\qquad
 a_n=3a_{n-1}-2a_{n-2}+a_{n-3}\quad(n\geq4).
\end{equation}
Equivalently, in the ring of formal power series,
\begin{equation}\label{eq:gf}
 \sum_{n\geq1}a_nz^n=\frac{z(1-z)^2}{1-3z+2z^2-z^3}.
\end{equation}
Moreover, for every \(n\geq1\) and \(s,t\in\B_n\),
\begin{equation}\label{eq:classification}
 s^A=t^A\quad\Longleftrightarrow\quad P(s)=P(t),
\end{equation}
where \(P\) is the label-preserving pruning map defined in
\eqref{eq:pruning} below.
\end{theorem}

The structural approach has precedents in the spectrum literature.
Cs\'ak\'any and Waldhauser~\cite[Section~4.2]{CsakanyWaldhauser2000}
use essential variables and recovery of subterms to separate bracketings
for the Sheffer operation. Huang and
Lehtonen~\cite[Section~7]{HuangLehtonen2023} study operations whose term
functions are classified by leaf-depth congruences. Here the complete
invariant is a pruned labeled comb, including the gaps in its leaf labels.
The contribution relative to the reported \(\SC\) data is the equality
classification and its all-arity enumeration.

\section{Pruning and reconstruction}

A tree is \emph{composite} if it is not a single leaf. A \emph{comb} is a
binary tree in which every internal node has a leaf child; a single leaf
also counts as a comb. Throughout, labels increase from left to right.
For any such labeled tree define
\begin{equation}\label{eq:pruning}
 P(x_i)=x_i,\qquad
 P(u*v)=
 \begin{cases}
  P(u),&u\text{ and }v\text{ are composite},\\
  P(u)*P(v),&\text{otherwise}.
 \end{cases}
\end{equation}
The test in this definition concerns the original children \(u,v\).
Surviving labels are never renumbered. A tree on a subset of the labels
defines a function on any common ambient set by ignoring the other variables.

\begin{lemma}\label{lem:pruning}
Pruning preserves the induced function. The tree \(P(t)\) is a comb,
retains the first leaf of \(t\), and has at least two leaves if \(t\)
is composite.
\end{lemma}
\begin{proof}
Every composite term has values in \(\{0,1\}\), and the restriction of
\(*\) to \(\{0,1\}\) is left projection. Thus \(u^A*v^A=u^A\)
whenever both children are composite. Structural induction now proves
preservation of the function in both cases of~\eqref{eq:pruning}.

The same induction gives the structural assertions. At a deleted root,
the pruned composite left child has at least two leaves and retains the
first leaf. At a retained root, at least one original child is a leaf;
the root therefore remains a comb node, and neither child is empty.
The two-leaf tree starts the composite case. In particular, pruning never
turns a composite child into a single leaf.
\end{proof}

A variable of a function is \emph{essential} if changing only that
coordinate can change the value.

\begin{lemma}\label{lem:essential}
For a comb with any increasing set of distinct labels, its leaf labels
are exactly the essential coordinates of its induced function.
\end{lemma}
\begin{proof}
The one-leaf case is immediate. For \(x_i*x_j\), varying \(x_i\)
between 0 and 1 at \(x_j=0\) witnesses the first coordinate;
varying \(x_j\) between 0 and 2 at \(x_i=1\) witnesses the second.

A comb with at least three leaves has exactly one composite child at
the root. If its function is \(g\), then \(g\) is Boolean-valued,
and the operation table gives the following three sections:
\begin{equation}\label{eq:sections}
\begin{aligned}
 \bigl((x_i*g)|_{x_i=0},(x_i*g)|_{x_i=1},(x_i*g)|_{x_i=2}\bigr)
   &=(0,1,1-g),\\
 \bigl((g*x_j)|_{x_j=0},(g*x_j)|_{x_j=1},(g*x_j)|_{x_j=2}\bigr)
   &=(g,g,0).
\end{aligned}
\end{equation}
Every essential coordinate of \(g\) remains essential, using \(x_i=2\)
in the first case and \(x_j=0\) in the second. The new left coordinate
is essential because its 0 and 1 sections differ. The new right coordinate
is essential because \(g=1\) when all its leaves equal 1, so its 0 and 2
sections differ. Absent coordinates have no effect. Induction completes
the proof.
\end{proof}


\begin{proposition}\label{prop:separation}
Two increasing-label combs induce the same function on a common ambient
set of variables if and only if they are identical as labeled trees.
\end{proposition}
\begin{proof}
Equality of functions gives equality of the leaf sets by
Lemma~\ref{lem:essential}. Induct on their common size \(k\). For
\(k=1,2\), the ordered tree on that set is unique. Suppose \(k\geq3\)
and let \(i,j\) be the least and greatest labels.

A root with a left leaf has the form \(x_i*u\). By
\eqref{eq:sections}, its \(x_i=1\) section is identically 1.
A root with a right leaf has the form \(u*x_j\), and its \(x_i=1\)
section is not constant: set every leaf of \(u\) to 1 and vary
\(x_j\) from 0 to 2. Thus the function determines the root type.
If \(f\) is the parent function, its composite child function is
recovered as
\[
 1-f|_{x_i=2}\quad\text{for a left leaf},\qquad
 f|_{x_j=0}\quad\text{for a right leaf}.
\]
The recovered functions agree on the remaining labels. Induction
determines the child trees, and hence the parent trees. The reverse
implication is immediate.
\end{proof}

Lemmas~\ref{lem:pruning} and~\ref{lem:essential} and
Proposition~\ref{prop:separation} prove~\eqref{eq:classification}.
The essential-coordinate set explains why gaps in the labels must be kept.
The reconstruction lemma applies even to combs that cannot arise by pruning
a bracketing on a consecutive interval.

There is a related reconstruction principle for weighted trees:
Zeng's Lemma~3.1~\cite{Zeng2026MultiplicityArxiv} recovers a weighted tree
with at least three vertices from weight multiplicities and counts of edges of each endpoint-weight
type when every nonleaf weight has multiplicity one. Here essential
coordinates identify the surviving labels, and the sections in
\eqref{eq:sections} successively recover the root and its child.
This comparison concerns recovery from an invariant; no weighted-tree
result is used in the proof.

\section{Enumeration of the attainable forms}

For \(n\geq2\), put
\[
 \Sset_n=\{P(t):t\in\B_n\},\qquad
 \Cset_n=\{p\in\Sset_n:\text{the largest label of }p\text{ is }n\},
 \qquad c_n=|\Cset_n|.
\]
By~\eqref{eq:classification}, \(a_n=|\Sset_n|\). Forms in different
sets \(\Cset_m\) have different largest labels. When comparing such
sets across arities, we retain their original labels.

\begin{lemma}\label{lem:terminal}
If \(p\in\Sset_n\) has largest label \(m\), then \(m\geq2\),
\(p\in\Cset_m\), and either \(m=n\) or \(m\leq n-2\).
Conversely, if \(p\in\Cset_m\) and \(n\geq m+2\), then
\(p\in\Sset_n\).
\end{lemma}
\begin{proof}
Lemma~\ref{lem:pruning} gives \(m\geq2\). Prove the forward
assertion by induction on the length of the label interval, allowing
uniform shifts of that interval. The case of two labels is immediate.
If \(m=n\), there is nothing to prove. Otherwise choose
\(t\in\B_n\) with \(P(t)=p\).

The root of \(t\) cannot have right child \(x_n\), which would
survive. If its left child is \(x_1\), its right child \(v\) is
composite on labels \(2,\ldots,n\). The shifted induction hypothesis
gives \(n-m\geq2\) and a bracketing \(v'\) on
\(2,\ldots,m\) with \(P(v')=P(v)\). Hence \(x_1*v'\) realizes
\(p\) on \(1,\ldots,m\).

In the remaining case both children are composite. Write the left
child \(u\) on \(1,\ldots,k\), where \(k\leq n-2\).
Then \(p=P(u)\) and \(m\leq k\). If \(m=k\), the child \(u\)
already realizes \(p\) on \(1,\ldots,m\); if \(m<k\), induction
provides such a realization. In both cases \(m\leq n-2\).
This proves \(p\in\Sset_m\), and therefore \(p\in\Cset_m\).

Conversely, take a realization \(t\in\B_m\) of \(p\) and any
bracketing \(v\) on the following \(n-m\geq2\) labels. Both
\(t\) and \(v\) are composite, so \(P(t*v)=P(t)=p\).
\end{proof}

Thus a deleted terminal interval cannot have length one. The lemma gives
the disjoint union and its counting identity
\begin{equation}\label{eq:decomposition}
 \Sset_n=\Cset_n\sqcup\bigsqcup_{m=2}^{n-2}\Cset_m,
 \qquad
 a_n=c_n+\sum_{m=2}^{n-2}c_m\qquad(n\geq2),
\end{equation}
with empty unions and sums understood in the usual way.


\begin{lemma}\label{lem:count}
We have \(c_2=1\) and
\begin{equation}\label{eq:c-rec}
 c_n=c_{n-1}+a_{n-1}\qquad(n\geq3).
\end{equation}
\end{lemma}
\begin{proof}
The unique element of \(\Cset_2\) is \(x_1*x_2\). A bracketing
realizing an element of \(\Cset_n\) cannot have two composite
children at its root, because that would delete \(x_n\).
For \(n\geq3\), there are two disjoint possibilities.

A left leaf \(x_1\) is joined to a pruned subtree on
\(2,\ldots,n\) retaining \(n\). These forms are in bijection with
\(\Cset_{n-1}\) by shifting all child labels by one. A right leaf
\(x_n\) is joined to any member of \(\Sset_{n-1}\), giving
\(a_{n-1}\) possibilities. Every form constructed this way has a
realization. The child opposite the leaf remains composite by
Lemma~\ref{lem:pruning}, so the two root types are disjoint. Each
construction is injective as a construction of labeled trees; multiple
realizations of a form are not counted separately.
\end{proof}

\begin{proof}[Completion of the proof of Theorem~\ref{thm:main}]
Let
\[
 C(z)=\sum_{n\geq2}c_nz^n,\qquad F(z)=\sum_{n\geq2}a_nz^n.
\]
Equations~\eqref{eq:decomposition} and~\eqref{eq:c-rec} give
\[
 F(z)=\left(1+\frac{z^2}{1-z}\right)C(z),\qquad
 C(z)=z^2+zC(z)+zF(z).
\]
Solving in formal power series yields
\[
 C(z)=\frac{z^2(1-z)}{1-3z+2z^2-z^3},\qquad
 F(z)=\frac{z^2(1-z+z^2)}{1-3z+2z^2-z^3}.
\]
Since the sole one-variable term is the identity, \(a_1=1\).
Adding \(z\) to \(F(z)\) proves~\eqref{eq:gf}. Multiplication by
its denominator gives \(z-2z^2+z^3\); comparison of coefficients
gives the three initial values and~\eqref{eq:recurrence} for every
\(n\geq4\). All denominators have constant term one. As a boundary
check, \(c_3=2\), \(c_4=4\), and \(a_4=c_4+c_2=5\).
\end{proof}

\section{Exact checks and scope}

The accompanying Python verifier enumerates full binary trees by splitting
each consecutive label interval at every possible root position.
It evaluates each tree on all ternary assignments, representing the positions
of values 1 and 2 by integer bit planes and applying all nine entries of
\eqref{eq:table}. Function equality is complete table equality. A separate
scalar table-lookup evaluator checks every tree and assignment through
\(n=7\). Tree generation and equality testing do not use the recurrence.

For every \(1\leq n\leq10\), the computation checks pruning preservation,
the equivalence~\eqref{eq:classification}, and equality of essential
support with the retained leaf set. It covers 6,918 bracketings in total,
including 4,862 bracketings and \(3^{10}=59,049\) assignments per bracketing
at \(n=10\). The resulting counts are
\[
\begin{array}{c|rrrrrrrrrr}
 n&1&2&3&4&5&6&7&8&9&10\\\hline
 a_n&1&1&2&5&12&28&65&151&351&816.
\end{array}
\]
The same program checks all 1,644 increasing-label combs on nonempty
subsets of eight labels, evaluating each on all \(3^8\) ambient
assignments. Their functions are distinct and their essential supports
are their leaf sets. This includes forms with gaps and forms not attainable
at that ambient arity. These are finite checks; the universal statements
follow from the preceding proofs.

The theorem concerns exactly the operation~\eqref{eq:table} and the
fixed-order spectrum. It supplies an equality test by pruning and counting
identities for every arity. The spectrum with variable permutations and
the other tables listed alongside \(\SC\) in
\cite[Section~7]{HuangLehtonen2024} are not treated here.

\subsection*{Computation and assistance}
The source package includes the standard-library Python verifier, its
recorded output, and reproduction instructions. OpenAI Codex assisted with
proof exploration, drafting, code generation, and a separate model-assisted
review of the argument. That review is not journal peer review. The checks
use exact integers, without sampling or floating-point decisions. No
proof-assistant formalization was performed.

\bibliographystyle{amsplain}
\bibliography{references}
\end{document}
